\paragraph{规范群熵账本、有限群代数障碍与偶 \texorpdfstring{$\zeta$}{zeta} 反演}
在规范群体积的二阶 Stirling 展开、审计偶数窗的最小退化律、约数触发零余类并的碰撞削减机制以及 fold-gauge 常数的 Stirling--Bernoulli 层级已经建立之后，还可进一步压缩出下列四条直接结论。

\begin{theorem}[规范群熵的精确账本恒等式]\label{thm:xi-time-part60aa-gauge-entropy-ledger-exact-identity}
沿用定理 \ref{thm:fold-bin-gauge-volume-stirling-second-order} 的记号，定义
\[
p_m^{\mathrm{bin}}(w):=\frac{d_m^{\mathrm{bin}}(w)}{2^m}
\qquad (w\in X_m),
\]
并记其 Shannon 熵为
\[
H\!\bigl(p_m^{\mathrm{bin}}\bigr):=
-\sum_{w\in X_m}p_m^{\mathrm{bin}}(w)\log p_m^{\mathrm{bin}}(w).
\]
则规范群体积密度 \(g_m\) 满足严格恒等式
\[
\boxed{\;
g_m
=
m\log 2
-H\!\bigl(p_m^{\mathrm{bin}}\bigr)
-1
+\frac{|X_m|\log(2\pi)+\sum_{w\in X_m}\log d_m^{\mathrm{bin}}(w)}{2^{m+1}}
+R_m,\;}
\]
其中余项满足
\[
\boxed{\;
0\le R_m\le \frac{|X_m|}{12\cdot 2^m}.\;}
\]
特别地，
\[
\boxed{\;
g_m
=
m\log 2
-H\!\bigl(p_m^{\mathrm{bin}}\bigr)
-1
+O\!\Bigl(\frac{|X_m|\,m}{2^m}\Bigr),\;}
\]
并且
\[
\boxed{\;
\kappa_m=m\log 2-H\!\bigl(p_m^{\mathrm{bin}}\bigr).\;}
\]
因此，\(\kappa_m\) 恰是 \(g_m\) 的零阶主项。
\end{theorem}
\begin{proof}
由定理 \ref{thm:fold-bin-gauge-volume-stirling-second-order}，
\[
g_m
=
\kappa_m-1+\frac{|X_m|\log(2\pi)+\sum_{w\in X_m}\log d_m^{\mathrm{bin}}(w)}{2^{m+1}}
+R_m,
\]
且
\[
0\le R_m\le \frac{|X_m|}{12\cdot 2^m}.
\]

另一方面，由
\[
p_m^{\mathrm{bin}}(w)=\frac{d_m^{\mathrm{bin}}(w)}{2^m},
\qquad
\sum_{w\in X_m}d_m^{\mathrm{bin}}(w)=2^m,
\]
可直接计算
\[
\begin{aligned}
H\!\bigl(p_m^{\mathrm{bin}}\bigr)
&=
-\sum_{w\in X_m}\frac{d_m^{\mathrm{bin}}(w)}{2^m}
\log\frac{d_m^{\mathrm{bin}}(w)}{2^m}\\
&=
-\frac{1}{2^m}\sum_{w\in X_m}d_m^{\mathrm{bin}}(w)\log d_m^{\mathrm{bin}}(w)
+\frac{m\log 2}{2^m}\sum_{w\in X_m}d_m^{\mathrm{bin}}(w)\\
&=
-\kappa_m+m\log 2.
\end{aligned}
\]
即
\[
\kappa_m=m\log 2-H\!\bigl(p_m^{\mathrm{bin}}\bigr).
\]
把该式代回 \(g_m\) 的 Stirling 展开，即得第一条盒装恒等式。

最后，由 \(1\le d_m^{\mathrm{bin}}(w)\le 2^m\) 可知
\[
0\le \sum_{w\in X_m}\log d_m^{\mathrm{bin}}(w)\le |X_m|\,m\log 2,
\]
再结合余项界，便得到
\[
g_m
=
m\log 2-H\!\bigl(p_m^{\mathrm{bin}}\bigr)-1
+O\!\Bigl(\frac{|X_m|\,m}{2^m}\Bigr).
\]
证毕。
\end{proof}

\begin{theorem}[审计偶数窗上的纤维群胚代数不可能是有限群代数]\label{thm:xi-time-part60aa-audited-even-groupoid-no-finite-group-algebra}
对审计偶数窗
\[
m\in\{6,8,10,12\},
\]
纤维群胚代数
\[
\CC[\mathcal R_m^{\mathrm{bin}}]
\cong
\bigoplus_{w\in X_m}M_{d_m^{\mathrm{bin}}(w)}(\CC)
\]
不可能与任何有限群 \(H\) 的复群代数 \(\CC[H]\) 同构。更强地，不存在任何非零代数同态
\[
\chi:\CC[\mathcal R_m^{\mathrm{bin}}]\to\CC.
\]
\end{theorem}
\begin{proof}
由命题 \ref{prop:fold-bin-groupoid-algebra-wedderburn}，
\[
\CC[\mathcal R_m^{\mathrm{bin}}]
\cong
\bigoplus_{w\in X_m}M_{d_m^{\mathrm{bin}}(w)}(\CC).
\]
而定理 \ref{thm:fold-bin-min-degeneracy-fib-audited} 在审计偶数窗上给出
\[
d_m^{\mathrm{bin}}(w)\ge d_{\min}(m)=F_{m/2}\ge 2
\qquad(\forall\,w\in X_m).
\]
因此上述 Wedderburn 分解中不存在任何 \(M_1(\CC)\) 块。

现证不存在非零代数同态 \(\chi:\CC[\mathcal R_m^{\mathrm{bin}}]\to\CC\)。若此类 \(\chi\) 存在，则其在某个直和因子
\[
M_{d_m^{\mathrm{bin}}(w)}(\CC)
\]
上的限制必非零。由于矩阵代数 \(M_n(\CC)\) 为简单代数，非零代数同态的核只能为零，故该限制必为单射。但当 \(n\ge 2\) 时，\(M_n(\CC)\) 非交换，而其像落在交换代数 \(\CC\) 中，矛盾。故不存在非零代数同态。

若反设存在有限群 \(H\) 使
\[
\CC[\mathcal R_m^{\mathrm{bin}}]\cong \CC[H],
\]
则群代数上的增广映射
\[
\varepsilon:\CC[H]\to\CC,
\qquad
\varepsilon\!\left(\sum_{h\in H}a_h h\right):=\sum_{h\in H}a_h
\]
给出一个非零代数同态。经由上述同构传回 \(\CC[\mathcal R_m^{\mathrm{bin}}]\)，便得到矛盾。故
\[
\CC[\mathcal R_m^{\mathrm{bin}}]\not\cong \CC[H]
\]
对一切有限群 \(H\) 都成立。证毕。
\end{proof}

\begin{theorem}[约数触发零余类并给出的规范群熵显式上界]\label{thm:xi-time-part60aa-divisor-zero-class-gauge-entropy-upper-bound}
记
\[
M:=F_{m+2}.
\]
若存在真约数 \(d\mid (m+2)\)、\(d<m+2\)，满足
\[
F_d\mid \frac{M}{2},
\]
则规范群体积密度 \(g_m\) 满足严格上界
\[
\boxed{\;
g_m
\le
m\log 2
-1
+\log\!\Bigl(1-\frac{F_d}{M}\Bigr)
+\frac{|X_m|\log(2\pi)+\sum_{w\in X_m}\log d_m^{\mathrm{bin}}(w)}{2^{m+1}}
+\frac{|X_m|}{12\cdot 2^m}.\;}
\]
特别地，
\[
\boxed{\;
g_m
\le
m\log 2
-1
+\log\!\Bigl(1-\frac{F_d}{M}\Bigr)
+O\!\Bigl(\frac{|X_m|\,m}{2^m}\Bigr).\;}
\]
\end{theorem}
\begin{proof}
由定理 \ref{thm:xi-time-part60aa-gauge-entropy-ledger-exact-identity}，
\[
g_m
=
m\log 2
-H\!\bigl(p_m^{\mathrm{bin}}\bigr)
-1
+\frac{|X_m|\log(2\pi)+\sum_{w\in X_m}\log d_m^{\mathrm{bin}}(w)}{2^{m+1}}
+R_m,
\]
其中
\[
0\le R_m\le \frac{|X_m|}{12\cdot 2^m}.
\]
另一方面，Shannon 熵支配 Rényi--\(2\) 熵，故
\[
H\!\bigl(p_m^{\mathrm{bin}}\bigr)\ge H_2\!\bigl(p_m^{\mathrm{bin}}\bigr)
=-\log \mathsf{Col}^{\mathrm{bin}}_m.
\]
于是
\[
g_m
\le
m\log 2
+\log \mathsf{Col}^{\mathrm{bin}}_m
-1
+\frac{|X_m|\log(2\pi)+\sum_{w\in X_m}\log d_m^{\mathrm{bin}}(w)}{2^{m+1}}
+\frac{|X_m|}{12\cdot 2^m}.
\]

再由推论 \ref{cor:fold-zero-divisor-lb}，
\[
|\mathcal Z_m|\ge F_d,
\]
而定理 \ref{thm:fold-collision-zero-reduction} 给出
\[
\mathsf{Col}^{\mathrm{bin}}_m
\le
1-\frac{|\mathcal Z_m|}{M}
\le
1-\frac{F_d}{M}.
\]
代回上式即得第一条盒装不等式。

最后，与定理 \ref{thm:xi-time-part60aa-gauge-entropy-ledger-exact-identity} 证明末尾相同，由
\[
0\le \sum_{w\in X_m}\log d_m^{\mathrm{bin}}(w)\le |X_m|\,m\log 2
\]
可将显示校正项统一吸收到
\(
O(|X_m|m/2^m)
\)
中，从而得到第二条公式。证毕。
\end{proof}

\begin{theorem}[从规范群常数逐阶显式恢复全部偶值 \texorpdfstring{$\zeta(2r)$}{zeta(2r)}]\label{thm:xi-time-part60aa-gauge-constant-explicit-even-zeta-recovery}
沿用定理 \ref{thm:fold-bin-gauge-constant-stirling-bernoulli-hierarchy} 的规范群常数列 \((C_m)_{m\ge 1}\)，并对每个整数 \(r\ge 1\) 定义
\[
A_r:=
\frac{B_{2r}}{r(2r-1)}
\bigl(\varphi^{-1}+\varphi^{2r-3}\bigr).
\]
再定义递阶剥离后的极限算子
\[
L_r:=
\lim_{m\to\infty}
\Bigl(\frac{2}{\varphi}\Bigr)^{(2r-1)m}
\left[
\log C_m-\log(2\pi)
-\sum_{s=1}^{r-1}A_s\Bigl(\frac{\varphi}{2}\Bigr)^{(2s-1)m}
\right].
\]
则极限存在且满足
\[
\boxed{\;L_r=A_r.\;}
\]
因此对每个 \(r\ge 1\)，全部偶值 \(\zeta(2r)\) 都可由显式公式恢复：
\[
\boxed{\;
\zeta(2r)
=
(-1)^{r-1}
\frac{(2\pi)^{2r}}{2(2r)!}\,
\frac{r(2r-1)}{\varphi^{-1}+\varphi^{2r-3}}\,
L_r.\;}
\]
\end{theorem}
\begin{proof}
由定理 \ref{thm:fold-bin-gauge-constant-stirling-bernoulli-hierarchy}，对任意固定 \(r\ge 1\) 有展开
\[
\log C_m
=
\log(2\pi)
+\sum_{s=1}^{r}A_s\Bigl(\frac{\varphi}{2}\Bigr)^{(2s-1)m}
+O\!\Bigl(\Bigl(\frac{\varphi}{2}\Bigr)^{(2r+1)m}\Bigr).
\]
将前 \(r-1\) 阶项移到左边，得到
\[
\log C_m-\log(2\pi)
-\sum_{s=1}^{r-1}A_s\Bigl(\frac{\varphi}{2}\Bigr)^{(2s-1)m}
=
A_r\Bigl(\frac{\varphi}{2}\Bigr)^{(2r-1)m}
+O\!\Bigl(\Bigl(\frac{\varphi}{2}\Bigr)^{(2r+1)m}\Bigr).
\]
两边同乘
\[
\Bigl(\frac{2}{\varphi}\Bigr)^{(2r-1)m},
\]
便得到
\[
\Bigl(\frac{2}{\varphi}\Bigr)^{(2r-1)m}
\left[
\log C_m-\log(2\pi)
-\sum_{s=1}^{r-1}A_s\Bigl(\frac{\varphi}{2}\Bigr)^{(2s-1)m}
\right]
=
A_r
+O\!\Bigl(\Bigl(\frac{\varphi}{2}\Bigr)^{2m}\Bigr).
\]
由于 \(\varphi/2<1\)，右端误差趋于 \(0\)，故极限存在且等于 \(A_r\)，即 \(L_r=A_r\)。

最后，由标准 Bernoulli--zeta 恒等式
\[
B_{2r}
=
(-1)^{r-1}\frac{2(2r)!}{(2\pi)^{2r}}\zeta(2r),
\]
以及
\[
A_r=
\frac{\varphi^{-1}+\varphi^{2r-3}}{r(2r-1)}\,B_{2r}
=L_r,
\]
解出 \(\zeta(2r)\) 即得所示恢复公式。证毕。
\end{proof}

\noindent
上述四条结论把规范群与 fold-gauge 常数的两条层级链收束为同一份显式账本：一方面，\(g_m\) 精确等于原始 \(m\)-比特预算减去折叠 Shannon 熵，并附带完全显示的 Stirling 校正；另一方面，\(\log C_m\) 的 Bernoulli 模态不仅可逐阶识别，而且可经由显式极限算子直接恢复全部偶值 \(\zeta(2r)\)。与审计偶数窗上的最小退化扇区及其零余类机制合并后，群熵损失、群代数实现障碍与偶 \(\zeta\)-塔反演遂被压缩进同一条闭合链。

\endinput
